\section{Background and Motivation}\label{sec:intro}

Every semester, University of Sydney students build their own timetable from the classes the university publishes~\cite{usyd_timetable}, choosing one class from every activity group so that nothing clashes. The combinations multiply quickly: ENGG2112 and AMME2301 alone allow \ttTwo{} clash-free timetables in Semester~2, 2026, even with full classes excluded, and five engineering units allow \ttFive{}.

For students who work, commute or care for family, the choice decides whether study fits around the rest of their life. In its 2017 Student Finances Survey, Universities Australia reported that more than four in five students were in paid work, 30\% of full-time students worked more than 20 hours a week, and more than a quarter of full-time students regularly missed classes because of work~\cite{ua_finances_2017}. These students state their requirements in words (``I work Tuesday and Thursday nights, and mornings aren't great''), not as weights or blocked grid cells.

Seen as an engineering problem, this is constrained scheduling with requirements written informally: some must hold (the student cannot attend), others are trade-offs (the student would rather not). UniPlan turns the sentence into explicit conditions, has the student confirm them, finds every timetable that satisfies the hard ones, and recommends three with a reason for each. The ``AI for Good'' case is that the students who gain most are the busiest, and that the system saves effort without deciding for them: every condition it reads is shown back, and the student can look past the top three.
